\pagestyle{empty}
\tight
\begin{tblr}{width=\textwidth,colspec={|Q[c,m,1.9cm]|X[l,m]|},hlines,vlines,hspan=minimal,row{1}={bg=headblue},rowsep=1pt,colsep=3pt}
\SetCell{c,m}\hfil\logo\hfil & \textbf{POLITEKNIK NEGERI MALANG}\par\textbf{JURUSAN TEKNOLOGI INFORMASI}\par\textbf{PROGRAM STUDI: D4 TEKNIK INFORMATIKA} \\
\SetCell[c=2]{c} \textbf{RENCANA PEMBELAJARAN SEMESTER (RPS)} & \\
\end{tblr}
\begin{longtblr}{width=\textwidth,colspec={|Q[l,m,1.9cm]|X[0.85,l,m]|X[1.45,l,m]|X[0.75,l,m]|X[1.15,l,m]|},hlines,vlines,hspan=minimal,rowsep=1pt,colsep=2pt,row{1,3}={bg=headgray,font=\bfseries}}
MATA KULIAH & KODE & BOBOT (SKS)/jam & SEMESTER & TGL. PENYUSUNAN \\
Sistem Informasi Manajemen & RTI253002 & 2 SKS/4 jam & 3 & 1 Juli 2025 \\
OTORISASI & Dosen Pengembang RPS & Koordinator RMK & \SetCell[c=2]{l} Ka PRODI &  \\
 & Agung Nugroho Pramudhita, S.T., M.T. & Agung Nugroho Pramudhita, S.T., M.T. & \SetCell[c=2]{l}{\vspace{8mm}\par Dr. Ely Setyo Astuti, S.T., M.T.} &  \\
\SetCell[r=13]{l} Capaian Pembelajaran (CP) & \SetCell[c=4]{l,bg=headgray,font=\bfseries} Capaian Pembelajaran Lulusan Program Studi (CPL-Prodi) &  &  &  \\
 & CPL02 & \SetCell[c=3]{l} Mampu menerapkan prinsip rekayasa sistem dan perangkat lunak untuk menghasilkan solusi aplikasi multi-platform sesuai kebutuhan. &  &  \\
 & CPL06 & \SetCell[c=3]{l} Mampu merancang, mengimplementasikan, menguji, melakukan deployment, memelihara, serta menjamin mutu perangkat lunak yang menjawab permasalahan dan memenuhi kebutuhan pengguna. &  &  \\
 & CPL08 & \SetCell[c=3]{l} Mampu mengelola proyek teknologi informasi secara profesional melalui perencanaan, koordinasi, komunikasi, dan optimalisasi sumber daya. &  &  \\
 & \SetCell[c=4]{l,bg=headgray,font=\bfseries} Capaian Pembelajaran mata kuliah (CPL-MK) &  &  &  \\
 & CPMK0211 & \SetCell[c=3]{l} Mahasiswa mampu menerapkan metodologi analisis dan desain sistem terstruktur menjadi spesifikasi teknis komprehensif untuk solusi sistem informasi. &  &  \\
 & CPMK0607 & \SetCell[c=3]{l} Mahasiswa mampu mengembangkan solusi sistem informasi manajemen terintegrasi yang kompleks sesuai kebutuhan organisasi. &  &  \\
 & CPMK0802 & \SetCell[c=3]{l} Mahasiswa mampu menerapkan rekayasa perangkat lunak dan prinsip SDLC dalam siklus hidup proyek sistem informasi. &  &  \\
 & \SetCell[c=4]{l,bg=headgray,font=\bfseries} Kemampuan akhir tiap tahapan belajar (Sub-CPMK) &  &  &  \\
 & SCPMK0211-02001 & \SetCell[c=3]{l} Mahasiswa mampu menganalisis konsep dasar sistem, informasi, dan struktur SIM untuk kebutuhan organisasi. &  &  \\
 & SCPMK0211-02002 & \SetCell[c=3]{l} Mahasiswa mampu merancang model konseptual SIM berdasarkan analisis proses bisnis dan alur keputusan. &  &  \\
 & SCPMK0607-02003 & \SetCell[c=3]{l} Mahasiswa mampu menerapkan konsep SI strategis dan pendukung untuk pengelolaan informasi organisasi yang kompleks. &  &  \\
 & SCPMK0802-02004 & \SetCell[c=3]{l} Mahasiswa mampu menyusun rencana pengembangan SIM menggunakan pendekatan SDLC yang tepat. &  &  \\
\SetCell[r=2]{l} Penilaian & \SetCell[c=4]{l} *Nilai CPMK didapat dari agregasi nilai Sub-CPMK yang dibebankan &  &  &  \\
 & \SetCell[c=4]{l}{\tiny\begin{tblr}{width=\linewidth,colspec={|Q[c,m,0.1\linewidth]|Q[c,m,0.16\linewidth]|X[l,m]|X[l,m]|X[l,m]|X[l,m]|X[l,m]|X[l,m]|Q[c,m,0.08\linewidth]|},hlines,vlines,hspan=minimal,rowsep=1pt,colsep=0.7pt,row{1,2}={bg=headgray,font=\bfseries}}
Kode CPMK & Kode SCPMK & \SetCell[c=6]{c} Bobot per Bentuk Penilaian & & & & & & Total Bobot \\
 & & Tugas & Kuis 1 & CM & Kuis 2 & --- & UAS & \\
\SetCell[r=2]{c} CPMK0211 & SCPMK0211-02001 & 10\%: konsep dasar sistem, informasi, dan struktur SIM & 5\%: konsep SIM & 0\% & 0\% & --- & 0\% & 15\% \\
 & SCPMK0211-02002 & 5\%: model konseptual dan proses bisnis & 0\% & 5\%: desain model SIM & 0\% & --- & 0\% & 10\% \\
\SetCell[r=1]{c} CPMK0607 & SCPMK0607-02003 & 5\%: SI strategis dan antar organisasi & 0\% & 30\%: studi kasus SI strategis & 5\%: SI pendukung & --- & 0\% & 40\% \\
\SetCell[r=1]{c} CPMK0802 & SCPMK0802-02004 & 0\% & 0\% & 20\%: rencana pengembangan SIM & 0\% & --- & 15\%: UAS SDLC & 35\% \\
\SetCell[c=8]{r}\textbf{Total Per Penilaian} & & & & & & & & \textbf{100\%} \\
\end{tblr}} &  &  &  \\
Diskripsi Singkat Mata Kuliah & \SetCell[c=4]{l} Mata kuliah Sistem Informasi Manajemen membangun kemampuan memahami konsep dasar, struktur, dan penerapan sistem informasi dalam organisasi, mencakup pengambilan keputusan berbasis informasi, SI strategis, dan pengembangan SIM berbasis SDLC. &  &  &  \\
Bahan Kajian / Materi Pembelajaran & \SetCell[c=4]{l}\begin{enumerate}\item Konsep dasar sistem dan informasi\item Struktur dan komponen SIM\item Organisasi dan pengambilan keputusan\item SI strategis dan antar organisasi\item Database, TIK, dan pengembangan SIM\end{enumerate} &  &  &  \\
\SetCell[r=2]{l} Pustaka & \SetCell[c=4]{l} \textbf{Utama:}\begin{enumerate}\item O'Brien, J.A. \& Marakas, G.M. 2011. Management Information Systems. McGraw-Hill.\item Laudon, K.C. \& Laudon, J.P. 2020. Management Information Systems. Pearson.\item McLeod, R. 2007. Sistem Informasi Manajemen. PT Indeks.\end{enumerate} &  &  &  \\
 & \SetCell[c=4]{l} \textbf{Pendukung:} Jurnal ilmiah, materi LMS, dan studi kasus aktual organisasi. &  &  &  \\
Media Pembelajaran & \SetCell[c=2]{l} \textbf{Software:} Office suite, diagram tools, LMS. &  & \SetCell[c=2]{l} \textbf{Hardware:} Komputer/laptop, proyektor. &  \\
Nama Dosen Pengampu & \SetCell[c=4]{l} Tim Pengajar Program Studi D4 TI &  &  &  \\
Matakuliah Syarat & \SetCell[c=4]{l} RTI252006 Basis Data &  &  &  \\
\end{longtblr}
\begin{longtblr}{width=\textwidth,colspec={|Q[c,m,0.06\textwidth]|X[1.35,l,m]|X[1.08,l,m]|X[1.16,l,m]|X[0.7,l,m]|X[1.24,l,m]|X[1.06,l,m]|X[0.98,l,m]|Q[c,m,0.05\textwidth]|},hlines,vlines,hspan=minimal,rowhead=2,row{1,2}={bg=headgreen,font=\bfseries},rowsep=1pt,colsep=1.5pt}
Mgg.\par Ke & Kemampuan Akhir Yang Direncanakan (Sub-CP-MK) & Bahan kajian (Materi Pembelajaran) & Bentuk dan Metode Pembelajaran & Estimasi Waktu & Pengalaman Belajar Mahasiswa & Kriteria \& Bentuk Penilaian & Indikator Penilaian & Bobot Penilaian (\%) \\
(1) & (2) & (3) & (4) & (5) & (6) & (7) & (8) & (9) \\
1 & SCPMK0211-02001: Mahasiswa mampu menganalisis konsep dasar sistem, informasi, dan struktur SIM untuk kebutuhan organisasi. & Konsep dasar sistem: pengertian sistem, karakteristik sistem, klasifikasi sistem, dan hubungan antar komponen dalam organisasi. & Modalitas: Blended Learning. Bentuk: Luring/Daring. Metode: Case Method, studi kasus, diskusi kelompok, dan presentasi. & 1 x 2 x 50' tatap muka; 1 x 2 x 50' tugas/praktik mandiri. & Mahasiswa mengerjakan aktivitas terarah untuk memahami konsep dasar sistem dan menunjukkan evidence hasil belajar. & Ketepatan konsep dasar sistem. Mengacu rubrik RTM. & Ketepatan konsep; kualitas analisis; validasi dan dokumentasi. & 3 \\
2 & SCPMK0211-02001: Mahasiswa mampu menganalisis konsep dasar sistem, informasi, dan struktur SIM untuk kebutuhan organisasi. & Konsep dasar dan perkembangan teknologi informasi: evolusi TI, komponen utama TI, peran TI dalam organisasi modern, dan tren perkembangan TI terkini. & Modalitas: Blended Learning. Bentuk: Luring/Daring. Metode: Case Method, studi kasus, diskusi kelompok, dan presentasi. & 1 x 2 x 50' tatap muka; 1 x 2 x 50' tugas/praktik mandiri. & Mahasiswa mengerjakan aktivitas terarah untuk memahami perkembangan teknologi informasi dan menunjukkan evidence hasil belajar. & Ketepatan konsep perkembangan TI. Mengacu rubrik RTM. & Ketepatan konsep; kualitas analisis; validasi dan dokumentasi. & 3 \\
3 & SCPMK0211-02001: Mahasiswa mampu menganalisis konsep dasar sistem, informasi, dan struktur SIM untuk kebutuhan organisasi. & Konsep dasar fakta, data, dan informasi: hierarki data, karakteristik informasi berkualitas, transformasi data menjadi informasi, dan nilai informasi dalam pengambilan keputusan. & Modalitas: Blended Learning. Bentuk: Luring/Daring. Metode: Case Method, studi kasus, diskusi kelompok, dan presentasi. & 1 x 2 x 50' tatap muka; 1 x 2 x 50' tugas/praktik mandiri. & Mahasiswa mengerjakan aktivitas terarah untuk memahami konsep fakta, data, dan informasi serta menunjukkan evidence hasil belajar. & Ketepatan konsep data dan informasi. Mengacu rubrik RTM. & Ketepatan konsep; kualitas analisis; validasi dan dokumentasi. & 3 \\
4 & SCPMK0211-02001: Mahasiswa mampu menganalisis konsep dasar sistem, informasi, dan struktur SIM untuk kebutuhan organisasi. & Konsep dan struktur sistem informasi manajemen: definisi SIM, komponen SIM, struktur hirarki SIM, hubungan SIM dengan proses bisnis, dan klasifikasi SIM berdasarkan fungsi manajemen. & Modalitas: Blended Learning. Bentuk: Luring/Daring. Metode: Case Method, studi kasus, diskusi kelompok, dan presentasi. & 1 x 2 x 50' tatap muka; 1 x 2 x 50' tugas/praktik mandiri. & Mahasiswa mengerjakan aktivitas terarah untuk memahami konsep dan struktur SIM serta menunjukkan evidence hasil belajar. & Ketepatan konsep dan struktur SIM. Mengacu rubrik RTM. & Ketepatan konsep; kualitas analisis; validasi dan dokumentasi. & 3 \\
5 & SCPMK0211-02001: Mahasiswa mampu menganalisis konsep dasar sistem, informasi, dan struktur SIM untuk kebutuhan organisasi. & Kuis 1: evaluasi pemahaman materi pertemuan 1--4 mencakup konsep sistem, perkembangan TI, fakta/data/informasi, dan struktur SIM. & Kuis tertulis (pilihan ganda dan/atau esai pendek). & 1 x 2 x 50'. & Mahasiswa menjawab soal kuis untuk menunjukkan pemahaman konsep dasar SIM. & Ketepatan menjawab soal. Bentuk: kuis tulis. & Ketepatan konsep; kualitas analisis; validasi dan dokumentasi. & 5 \\
6 & SCPMK0211-02002: Mahasiswa mampu merancang model konseptual SIM berdasarkan analisis proses bisnis dan alur keputusan. & Konsep organisasi dalam SIM: struktur organisasi, level manajemen, kebutuhan informasi tiap level, hubungan SIM dengan fungsi organisasi, dan model organisasi berbasis informasi. & Modalitas: Blended Learning. Bentuk: Luring/Daring. Metode: Case Method, studi kasus, diskusi kelompok, dan presentasi. & 1 x 2 x 50' tatap muka; 1 x 2 x 50' tugas/praktik mandiri. & Mahasiswa mengerjakan aktivitas terarah untuk memahami konsep organisasi dalam SIM dan menunjukkan evidence hasil belajar. & Ketepatan konsep organisasi dalam SIM. Mengacu rubrik RTM. & Ketepatan konsep; kualitas analisis; validasi dan dokumentasi. & 3 \\
7 & SCPMK0211-02002: Mahasiswa mampu merancang model konseptual SIM berdasarkan analisis proses bisnis dan alur keputusan. & Pengambilan keputusan berbasis SI: model pengambilan keputusan, peran SI dalam mendukung keputusan, DSS dan MIS, alur pengambilan keputusan pada berbagai level manajemen, dan studi kasus. & Modalitas: Blended Learning. Bentuk: Luring/Daring. Metode: Case Method, studi kasus, diskusi kelompok, dan presentasi. & 1 x 2 x 50' tatap muka; 1 x 2 x 50' tugas/praktik mandiri. & Mahasiswa mengerjakan aktivitas terarah untuk memahami pengambilan keputusan berbasis SI dan menunjukkan evidence hasil belajar. & Ketepatan konsep pengambilan keputusan berbasis SI. Mengacu rubrik RTM. & Ketepatan konsep; kualitas analisis; validasi dan dokumentasi. & 3 \\
8 & SCPMK0211-02001, SCPMK0211-02002: Mahasiswa mampu menganalisis konsep dasar SIM dan merancang model konseptual SIM berdasarkan analisis proses bisnis. & UTS: materi pertemuan 1--7 mencakup konsep sistem, TI, data/informasi, struktur SIM, organisasi, dan pengambilan keputusan berbasis SI. & Ujian tertulis berbasis analisis kasus. & 1 x 2 x 50'. & Mahasiswa menganalisis kasus organisasi dan merancang model SIM secara tertulis untuk menunjukkan evidence hasil belajar. & Ketepatan analisis dan pemodelan SIM. Bentuk: tes tulis. & Ketepatan analisis kebutuhan organisasi dan perancangan model konseptual SIM. & 15 \\
9 & SCPMK0607-02003: Mahasiswa mampu menerapkan konsep SI strategis dan pendukung untuk pengelolaan informasi organisasi yang kompleks. & Konsep teknologi dan SI untuk SIM: infrastruktur TI, jaringan komputer, cloud computing, peran teknologi dalam mendukung SIM, dan integrasi teknologi dalam sistem informasi organisasi. & Modalitas: Blended Learning. Bentuk: Luring/Daring. Metode: Case Method, studi kasus, diskusi kelompok, dan presentasi. & 1 x 2 x 50' tatap muka; 1 x 2 x 50' tugas/praktik mandiri. & Mahasiswa mengerjakan aktivitas terarah untuk memahami konsep teknologi SI dan menunjukkan evidence hasil belajar. & Ketepatan konsep teknologi SI untuk SIM. Mengacu rubrik RTM. & Ketepatan konsep; kualitas analisis; validasi dan dokumentasi. & 5 \\
10 & SCPMK0607-02003: Mahasiswa mampu menerapkan konsep SI strategis dan pendukung untuk pengelolaan informasi organisasi yang kompleks. & Aplikasi SI pada fungsi dan level organisasi: SI fungsional (keuangan, SDM, pemasaran, operasi), SI pada level operasional/taktis/strategis, ERP, CRM, SCM, dan studi kasus penerapan SI. & Modalitas: Blended Learning. Bentuk: Luring/Daring. Metode: Case Method, studi kasus, diskusi kelompok, dan presentasi. & 1 x 2 x 50' tatap muka; 1 x 2 x 50' tugas/praktik mandiri. & Mahasiswa mengerjakan aktivitas terarah untuk menganalisis aplikasi SI pada organisasi dan menunjukkan evidence hasil belajar. & Ketepatan analisis aplikasi SI pada organisasi. Mengacu rubrik RTM. & Ketepatan konsep; kualitas analisis; validasi dan dokumentasi. & 5 \\
11 & SCPMK0607-02003: Mahasiswa mampu menerapkan konsep SI strategis dan pendukung untuk pengelolaan informasi organisasi yang kompleks. & Konsep SI strategis dan sistem antar organisasi: SI strategis untuk keunggulan kompetitif, Value Chain, Porter's Five Forces, sistem antar organisasi (IOS), EDI, dan ekosistem digital bisnis. & Modalitas: Blended Learning. Bentuk: Luring/Daring. Metode: Case Method, studi kasus, diskusi kelompok, dan presentasi. & 1 x 2 x 50' tatap muka; 1 x 2 x 50' tugas/praktik mandiri. & Mahasiswa mengerjakan aktivitas terarah untuk menganalisis SI strategis dan antar organisasi serta menunjukkan evidence hasil belajar. & Ketepatan analisis SI strategis. Mengacu rubrik RTM. & Ketepatan konsep; kualitas analisis; validasi dan dokumentasi. & 5 \\
12 & SCPMK0607-02003: Mahasiswa mampu menerapkan konsep SI strategis dan pendukung untuk pengelolaan informasi organisasi yang kompleks. & SI pendukung dalam SIM: Decision Support System (DSS), Executive Information System (EIS), Expert System, Knowledge Management System (KMS), Business Intelligence (BI), dan peran SI pendukung dalam pengambilan keputusan manajerial. & Modalitas: Blended Learning. Bentuk: Luring/Daring. Metode: Case Method, studi kasus, diskusi kelompok, dan presentasi. & 1 x 2 x 50' tatap muka; 1 x 2 x 50' tugas/praktik mandiri. & Mahasiswa mengerjakan aktivitas terarah untuk memahami SI pendukung dalam SIM dan menunjukkan evidence hasil belajar. & Ketepatan konsep SI pendukung dalam SIM. Mengacu rubrik RTM. & Ketepatan konsep; kualitas analisis; validasi dan dokumentasi. & 5 \\
13 & SCPMK0607-02003: Mahasiswa mampu menerapkan konsep SI strategis dan pendukung untuk pengelolaan informasi organisasi yang kompleks. & Kuis 2: evaluasi pemahaman materi pertemuan 9--12 mencakup konsep teknologi SI, aplikasi SI organisasi, SI strategis, dan SI pendukung. & Kuis tertulis (pilihan ganda dan/atau esai pendek). & 1 x 2 x 50'. & Mahasiswa menjawab soal kuis untuk menunjukkan pemahaman konsep SI strategis dan pendukung. & Ketepatan menjawab soal. Bentuk: kuis tulis. & Ketepatan konsep; kualitas analisis; validasi dan dokumentasi. & 5 \\
14 & SCPMK0802-02004: Mahasiswa mampu menyusun rencana pengembangan SIM menggunakan pendekatan SDLC yang tepat. & Basis data dan DBMS: konsep basis data relasional, peran DBMS dalam SIM, desain basis data untuk SIM, integrasi data dalam sistem informasi, dan pengelolaan data organisasi. & Modalitas: Blended Learning. Bentuk: Luring/Daring. Metode: Case Method, studi kasus, diskusi kelompok, dan presentasi. & 1 x 2 x 50' tatap muka; 1 x 2 x 50' tugas/praktik mandiri. & Mahasiswa mengerjakan aktivitas terarah untuk memahami peran basis data dan DBMS dalam SIM serta menunjukkan evidence hasil belajar. & Ketepatan konsep basis data dan DBMS untuk SIM. Mengacu rubrik RTM. & Ketepatan konsep; kualitas analisis; validasi dan dokumentasi. & 5 \\
15 & SCPMK0802-02004: Mahasiswa mampu menyusun rencana pengembangan SIM menggunakan pendekatan SDLC yang tepat. & Teknologi Informasi dan Komunikasi (TIK) untuk SIM: infrastruktur TIK, jaringan dan komunikasi data, keamanan sistem informasi, TIK pendukung SIM modern, dan perencanaan TIK dalam proyek SIM. & Modalitas: Blended Learning. Bentuk: Luring/Daring. Metode: Case Method, studi kasus, diskusi kelompok, dan presentasi. & 1 x 2 x 50' tatap muka; 1 x 2 x 50' tugas/praktik mandiri. & Mahasiswa mengerjakan aktivitas terarah untuk memahami peran TIK dalam SIM dan menunjukkan evidence hasil belajar. & Ketepatan konsep TIK untuk SIM. Mengacu rubrik RTM. & Ketepatan konsep; kualitas analisis; validasi dan dokumentasi. & 5 \\
16 & SCPMK0802-02004: Mahasiswa mampu menyusun rencana pengembangan SIM menggunakan pendekatan SDLC yang tepat. & Pengembangan SIM: metodologi SDLC (Waterfall, Agile, Spiral), tahapan pengembangan sistem, analisis kebutuhan, desain sistem, implementasi, pengujian, dan pemeliharaan SIM. & Modalitas: Blended Learning. Bentuk: Luring/Daring. Metode: Case Method, studi kasus, diskusi kelompok, dan presentasi. & 1 x 2 x 50' tatap muka; 1 x 2 x 50' tugas/praktik mandiri. & Mahasiswa mengerjakan aktivitas terarah untuk merancang rencana pengembangan SIM menggunakan SDLC dan menunjukkan evidence hasil belajar. & Ketepatan rencana pengembangan SIM berbasis SDLC. Mengacu rubrik RTM. & Ketepatan konsep; kualitas analisis; validasi dan dokumentasi. & 5 \\
17 & SCPMK0607-02003, SCPMK0802-02004: Mahasiswa mampu menerapkan SI strategis dan menyusun rencana pengembangan SIM menggunakan pendekatan SDLC yang tepat. & UAS: materi pertemuan 9--16 dengan fokus pada penerapan SI strategis dan penyusunan rencana pengembangan SIM menggunakan SDLC. & UAS berbasis proyek mini. & 1 x 2 x 50'. & Mahasiswa menyusun rencana pengembangan SIM berbasis SDLC untuk konteks organisasi yang diberikan dan menunjukkan evidence hasil belajar. & Ketepatan pendekatan SDLC dan kelengkapan rencana pengembangan SIM. Bentuk: proyek mini. & Ketepatan pendekatan SDLC; kelengkapan rencana pengembangan; kualitas spesifikasi teknis. & 22 \\
\end{longtblr}
